\documentclass[11pt]{article}
\usepackage[T1]{fontenc}
\usepackage[utf8]{inputenc}
\usepackage{lmodern}
\usepackage[ngerman]{babel}
\usepackage{amsmath,amssymb,mathtools,array,booktabs}
\usepackage{geometry}
\usepackage{microtype}
\geometry{a4paper,margin=1.45cm}
\setlength{\parindent}{1.25em}
\setlength{\parskip}{0pt}
\makeatletter
\renewcommand\footnotesize{\@setfontsize\footnotesize{7.5}{8.5}\selectfont}
\makeatother
\allowdisplaybreaks
\setlength{\emergencystretch}{10em}
\sloppy
\hbadness=10000
\vbadness=10000
\hfuzz=2pt
\vfuzz=10pt
\raggedbottom
\providecommand{\p}{\partial}
\providecommand{\dd}{\mathrm d}
\providecommand{\modu}[1]{\;\operatorname{mod} #1}
\begin{document}

\section*{16. Zur Reihenentwicklung in der Formentheorie}
\begin{center}
Math. Ann. 81 (1920), S. 25--30
\end{center}

Unter ,,Reihenentwicklung der Formentheorie`` versteht man bekanntlich einerseits die Verallgemeinerung der Clebsch-Gordanschen Reihenentwicklung -- d. h. der Entwicklung einer zwei binäre Reihen enthaltenden Form nach Potenzen der Determinante dieser Reihen, wobei die Koeffizienten Polaren von Formen von nur einer Variabelnreihe sind, oder, was damit identisch ist, bei Anwendung des $\Omega$-Prozesses verschwinden -- auf Formen von beliebig vielen Serien von je $n$ Variabeln; anderseits als Verallgemeinerung der bei ,,Komplexkoordinaten`` $t_{ik}$ auftretenden Normalformen eine Entwicklung nach Normalformen für Formen, die die Größen verschiedener Stufen $t_i,t_{ik},t_{ikl},\ldots$ gleichzeitig enthalten. Diese zweite Art der Entwicklung läßt sich aber aus der ersten durch invariante Differentiationsprozesse gewinnen, wie dies vor allem Mertens und Deruyts ausgeführt haben.\footnote{F. Mertens, Über invariante Gebilde quaternärer Formen (Sitzber. d. Ak. d. Wiss. Wien 98, Abt. IIa (1889)) für den quaternären Fall. Die allgemeinen Normalformen finden sich bei J. Deruyts: Sur la réduction des fonctions invariantes dans le système des variables géométriques, Bulletins de l'Ac. R. de Belgique 24 (1892). Ohne Kenntnis dieser Arbeit habe ich in genauer Anlehnung an Mertens die allgemeine Entwicklung nach Normalformen angegeben in \S~7 der Arbeit: Zur Invariantentheorie der Formen von $n$ Variabeln, J. f. M. 139 (1910).}

Die verschiedenen Entwicklungen erster Art ergeben sich alle -- wie dies etwa Math. Ann. 77 a. a. O. III kurz skizziert -- als Folgerung der Entwicklung einer Form von $n$ Reihen von je $n$ Variabeln nach Potenzen der Determinante dieser Reihen, wobei die Koeffizienten Polaren von Formen von nur $(n-1)$ Reihen sind, die selbst aus der Ausgangsform durch Polarenbildung und $\Omega$-Prozeß abgeleitet sind.\footnote{Vgl. die Literaturangaben a. a. O. Seitdem erschienen und auch auf der formalen Methode beruhend: Schouten, Over reeksontwikkelingen van algebraische vormen met verschillende rijen van variabelen van verschillenden graad, Verslag Ak. Amsterdam 27 (1919), S. 1481.} Die ursprüngliche Clebsch-Gordansche Entwicklung $(n=2)$ war noch einen Schritt weiter gegangen: die speziellen auftretenden Polarprozesse sind dort explizite angegeben, sogar unter Einschluß der numerischen Koeffizienten.

Was ich nun der angegebenen Note -- wo die Reihenentwicklung mehr begrifflich, als ein Problem der linearen Formenscharen aufgefaßt ist -- zufügen möchte, ist eine bis auf Zahlenkoeffizienten explizite Darstellung, die sich durch Spezialisierung der auf allgemeine Modulsysteme bezüglichen Untersuchungen von E. Fischer\footnote{E. Fischer, Über die Differentiationsprozesse der Algebra, J. f. M. 148 (1917).} ergibt.

Zu dieser Einordnung führt die Deutung der Normalform in den $t_{ik}$. Ersetzt man nämlich die ,,Komplexkoordinaten`` $t_{ik}$ durch die ,,Linienkoordinaten``
\[
        p_{ik}=(x_i y_k-x_k y_i),
\]
so ist vermöge der zwischen den $p_{ik}$ bestehenden identischen Relationen die Darstellung einer Form $F(p_{ik})$ nicht mehr eindeutig. Eindeutigkeit läßt sich erzielen, wenn man verlangt, daß zwischen den Koeffizienten von $F$ dieselben linearen Relationen bestehen wie zwischen den in $F$ auftretenden Potenzprodukten der $p_{ik}$.\footnote{Vgl. Clebsch, Über eine Fundamentalaufgabe der Invariantentheorie, \S~4, Abh. d. K. G. d. Wissensch. zu Göttingen 17 (1872).} Für diese Normalform $\Phi$ gilt also:
\begin{equation}
\begin{aligned}
        F(p_{ik})&=\Phi(p_{ik}) &&[\text{identisch in }x,y],\quad\text{oder}\notag\\
        F(t_{ik})&\equiv\Phi(t_{ik}) &&\modu{M},
\end{aligned}
\tag{1}
\end{equation}
wo $M$ den Modul aller identischen Relationen zwischen den $p_{ik}$ bedeutet, d. h. aus der Gesamtheit aller Formen $G(t_{ik})$ besteht, für die $G(p_{ik})$ identisch in $x,y$ verschwindet. $\Phi(t_{ik})$ ist die ,,Normalform`` in Komplexkoordinaten, während die Basisfunktionen von $M$ quadratische Formen in den $t_{ik}$ werden. Durch Iteration des Verfahrens in bezug auf die Faktoren dieser Basisfunktionen entsteht die vollständige Reihenentwicklung. Entsprechende Überlegungen gelten bei gleichzeitigem Auftreten der $t_i,t_{ik},t_{ikl},\ldots$.

Auf einen analogen Gedankengang läßt sich die Entwicklung nach Polaren zurückführen. Ersetzt man nämlich etwa in einer Form $F(x,y,z)$ die drei ternären Reihen $x=(x_1,x_2,x_3),y,z$ durch solche Reihen $\xi,\eta,\zeta$, die sich linear aus nur zwei zusammensetzen lassen (etwa $\xi=x,\eta=y;\zeta=\lambda_1x+\lambda_2y$), so wird die Darstellung von $F(\xi,\eta,\zeta)$ wieder nicht eindeutig. Man kann wieder für die Koeffizienten dieselben linearen Relationen vorschreiben, wie für die in $F$ auftretenden Potenzprodukte der $\xi,\eta,\zeta$. Da nun hier der Modul $M$ aller identischen Relationen als Basisfunktion die Determinante $\Delta=(xyz)$ der drei Reihen besitzt (a. a. O. III, Hilfssatz a)), so tritt hier an Stelle von [1]:
\begin{equation}
\begin{aligned}
        F(\xi,\eta,\zeta)&=\Phi(\xi,\eta,\zeta)&&[\text{identisch in }x,y],\quad\text{oder}\notag\\
        F(x,y,z)&\equiv\Phi(x,y,z)&&\modu{\Delta}.
\end{aligned}
\tag{2}
\end{equation}
Es zeigt sich, daß $\Phi$ durch Polarenbildung entsteht (vgl. Formel (2) a. a. O.); die vollständige Reihenentwicklung ergibt sich wieder durch Iteration.

Die Bildung der Normalformen $\Phi$ ordnet sich also beidemal folgendem Problem unter: Ersetzt man in $F(z)$ die Unbestimmten $z_i$ durch Polynome $f_i(\xi)$ von Parametern $\xi$, so daß also die Darstellung $F(f_i(\xi))$ nicht mehr eindeutig wird, so soll die Eindeutigkeit dadurch festgelegt werden, daß zwischen den Koeffizienten dieselben linearen Relationen bestehen wie zwischen den in $F$ auftretenden Potenzprodukten der $f_i(\xi)$. Es wird also
\begin{equation}
\begin{aligned}
        F(f_i(\xi))&=\Phi(f_i(\xi))&&[\text{identisch in }\xi],\quad\text{oder}\notag\\
        F(z)&\equiv\Phi(z)&&\modu{M},
\end{aligned}
\tag{3}
\end{equation}
wo $M$ den Modul aller identischen Relationen zwischen den $f_i(\xi)$ bedeutet.

Für diese Normalform $\Phi(z)$ hat nun E. Fischer in \S~11, I seiner zitierten Arbeit eine bis auf Zahlkoeffizienten explizite Entwicklung, und zwar vermöge linearer Operationen angegeben. Zugleich hat er aber auch in der Einleitung III und in \S~6 II für die zu $M$ gehörige Differenz $F-\Phi$ -- unter Voraussetzung der Kenntnis einer Modulbasis -- einen im selben Sinn expliziten Ausdruck angegeben; und in Verbindung dieser Entwicklungen in \S~11 II die Formel [3] somit durch eine explizite Darstellung ersetzt. Dabei ist die noch offene Frage nach den Zahlkoeffizienten nach \S~12 derselben Arbeit identisch mit dem Aufsuchen der Eigenwerte und Eigenfunktionen der linearen Operationen (Formel 82).

Ich spezialisiere nun in 1. die Fischerschen Formeln auf den Fall der Reihenentwicklung nach Polaren und gewinne daraus durch Iteration die Reihenentwicklung. Die Operation zur Bestimmung von $\Phi$ geht dabei in bestimmte Polarprozesse über, was erheblich genauer ist als die oben angegebene, bisher bekannte Fassung. Bei der Operation zur Bestimmung von $F-\Phi$ tritt die Multiplikation mit der Determinante $\Delta$ und der $\Omega$-Prozeß gemischt auf. Um auf die gewöhnliche, nicht explizite Form zu kommen, ist noch die Bemerkung zu machen, daß $\Omega\Delta$ sich durch Polarprozesse ausdrücken läßt.

An dieser Stelle sei noch eine Berichtigung zugefügt. Ich habe a. a. O. (S. 101, Z. 6 v. o.) versehentlich die nur im binären Gebiet geltende identische Umformung
\[
        \Omega(\Delta\psi)=c_1\psi+c_2\Delta\cdot\Omega(\psi)
\]
als allgemeingültig angenommen und damit den Übergang von der Fundamentalformel (2), die den Ausdruck für die Normalform $\Phi$ in [2] gibt, auf die Reihenentwicklung (3) gemacht. Die dortigen Überlegungen ergeben also nur die Formel (2), und als Spezialfall davon den Reduktionssatz (1); reichen aber, da ein expliziter Ausdruck für die Differenz $F-\Phi$ fehlt, zur Begründung der Reihenentwicklung nicht aus.

In 2. erhalte ich ebenso durch Spezialisierung die Normalform $\Phi(t_{ik})$ von [1]; allgemeiner die Normalform $\Phi(t_i,t_{ik},t_{ikl},\ldots)$, was auf bekannte explizite Formeln führt. Da aber die Aufstellung der vollständigen Reihenentwicklung die Kenntnis einer Modulbasis von $M$ verlangt, ist hier der invariantentheoretische Weg (vgl. die angeführte Literatur) vorzuziehen, der diese Kenntnis nicht verlangt, sondern im Gegenteil vermöge der Reihenentwicklung eine Modulbasis liefert.

1. Nach Fischer, \S~11 (Formel (19) und folgende Zeilen) wird die Normalform $\Phi(z)$ von [3] -- (dort $N(\zeta)$) -- gegeben durch:
\begin{equation}
        \Phi(z)=(a_1S+a_2S^2+\cdots+a_rS^r)F(z),
\tag{4}
\end{equation}
wo die Konstanten $a_1\ldots a_r$ dem Koeffizientenbereich der $f(\xi)$ angehören und nur vom Grade von $F(z)$ abhängen; der Operator $S$ ist definiert durch
\begin{equation}
\begin{gathered}
        S=AB;\qquad BF(z)=F(f(\xi));\\
        AF(f(\xi))=\left\{\sum z_i f_i\!\left(\frac{\p}{\p\xi}\right)\right\}^{p}F(f(\xi))\footnotemark.
\end{gathered}
\tag{5}
\end{equation}
\footnotetext{Vgl. Formel (75). Im Falle von [3] läuft die Summe $\alpha$ über alle Potenzprodukte $p$-ter Dimensionen, und (75) spezialisiert sich daher zu [5].}
Faßt man in [2] $F$ als Form $p$-ter Dimension in $z$ allein auf, so spezialisiert sich $f_i(\xi)$ zu $\lambda_1x_i+\lambda_2y_i$; somit kommt:
\[
\begin{gathered}
        BF=F(x,y,\lambda_1x+\lambda_2y),\\
        SF=\left\{\sum z_i\left(\frac{\p}{\p\lambda_1}\frac{\p}{\p x_i}+
        \frac{\p}{\p\lambda_2}\frac{\p}{\p y_i}\right)\right\}^{p}
        F(x,y,\lambda_1x+\lambda_2y),
\end{gathered}
\]
oder in Polarprozessen geschrieben (vgl. a. a. O. II, 2. u. 3.; auch für die abkürzende Bezeichnung):
\begin{equation}
        SF(x,y,z)=\frac1{p!}\left[z\left(\frac{\p}{\p\lambda_1}\frac{\p}{\p x}+
        \frac{\p}{\p\lambda_2}\frac{\p}{\p y}\right)\right]^p
        \left[(\lambda_1x+\lambda_2y)\frac{\p}{\p z}\right]^p F(x,y,z).
\tag{6}
\end{equation}
\emph{Durch} [4] \emph{und} [6] \emph{ist der gewünschte explizite Ausdruck für} $\Phi(x,y,z)$ \emph{gegeben}; $\Phi(x,y,z)$ ordnet sich zugleich der am Anfang erwähnten allgemeinen Form ($\sum PH$ in (2), a. a. O.) -- Polaren von Ausdrücken, die nur von 2 Reihen abhängen und aus $F$ durch Polarprozesse gewonnen werden -- unter. Die noch unbestimmten Koeffizienten $a_i$ sind rationale Zahlen; da hier die $f_i(\xi)$ ganzrationale Zahlkoeffizienten besitzen. Schreibt man [2] in der Form:
\begin{equation}
        F=\Phi+\Delta F_1,
\tag{2a}
\end{equation}
so kommt durch Spezialisierung der bei Fischer in der Einleitung III und in \S~6 II (S. 41) angegebenen Formel
\begin{equation}
        F_1=(c_1\Omega+c_2\Omega\Delta\Omega+\cdots+c_i\Omega\Delta\Omega\cdots\Delta\Omega)F\footnotemark,
        \qquad \left[\Omega=\left(\frac{\p}{\p x}\frac{\p}{\p y}\frac{\p}{\p z}\right)\right],
\tag{7}
\end{equation}
\footnotetext{Diese spezialisierte Formel benutzt Kerim in seiner demnächst erscheinenden Erlanger Dissertation in der Anwendung auf ,,Trägheitsformen``.}
wo wieder die $c_i$ rationale, nur vom Grad von $F$ abhängende Zahlen sind.

Die Iteration führt zu:
\[
        F_1=\Phi_1+\Delta F_2;\quad F_2=\Phi_2+\Delta F_3;\quad\ldots;\quad
        F_i=\Phi_i+\Delta F_{i+1};\quad\ldots,
\]
wo jeweils die $\Phi_i$ die den $F_i$ entsprechenden Normalformen sind. Da die $F_i$ in $z$ vom Grade $(p-i)$ sind, kommt nach [4] und [6]:
\begin{equation}
\begin{gathered}
        \Phi_i=(a_{i1}S_i+a_{i2}S_i^2+\cdots)F_i,\\
        S_iF_i=\frac1{(p-i)!}
        \left[z\left(\frac{\p}{\p\lambda_1}\frac{\p}{\p x}+
        \frac{\p}{\p\lambda_2}\frac{\p}{\p y}\right)\right]^{p-i}
        \left[(\lambda_1x+\lambda_2y)\frac{\p}{\p z}\right]^{p-i}F_i.
\end{gathered}
\tag{8}
\end{equation}
Der Formel [7] entspricht:
\[
        F_i=(\alpha_{i1}\Omega+\alpha_{i2}\Omega\Delta\Omega+\cdots)F_{i-1},
        \quad\text{und durch Iteration}
\]
\begin{equation}
        F_i=(\beta_1\Omega^i+\beta_2\Omega^i\Delta\Omega+\cdots
        +\beta_\rho\Omega^{k_1}\Delta\Omega^{k_2}\cdots\Delta\Omega^{k_\lambda}+\cdots)F,
        \quad(k_1+k_2+\cdots+k_\lambda-\lambda+1=i).
\tag{9}
\end{equation}
\emph{Vermöge} [8] \emph{und} [9] \emph{ist in der Reihenentwicklung}
\[
        F=\Phi+\Delta\Phi_1+\Delta^2\Phi_2+\cdots+\Delta^i\Phi_i+\cdots+\Delta^p\Phi_p
\]
\emph{die im angegebenen Sinn explizite Darstellung erreicht.}

Wendet man auf [7] die identische Umformung: $\Omega\Delta F=PF$ an, wo $P$ Polarprozesse bedeuten,\footnote{Vgl. etwa F. Mertens, Über ganze Funktionen von $m$ Systemen von je $n$ Unbestimmten, Monatsh. f. Math. u. Phys. 4, Formel (5).} so kommt, wegen der Vertauschbarkeit dieser Prozesse mit dem $\Omega$-Prozeß, die bekannte Formel:
\[
        F_1=P\Omega F;\quad\ldots;\quad F_i=P\Omega^iF;
\]
und daraus folgt in Verbindung mit [8] oder mit den Überlegungen a. a. O. die bekannte Fassung der Reihenentwicklung, wie etwa in (3) a. a. O. Handelt es sich um $n$ Reihen von $n$ Variabeln, $x^{(1)},x^{(2)},\ldots,x^{(n-1)},z$; so wird entsprechend:
\[
        f_i(\xi)=\lambda_1x_i^{(1)}+\lambda_2x_i^{(2)}+\cdots+\lambda_{n-1}x_i^{(n-1)};
\]
und ebenso ist $\Delta$ und $\Omega$ für $n$ Reihen zu definieren.

2. Für die Normalform der $F(t_i,t_{ik},t_{ikl},\ldots)$ sind die $f_i(\xi),f_{ik}(\xi),\ldots$ gegeben durch die ein-, zwei- bis $(n-1)$-reihigen Determinanten
\[
        \xi_i^{(1)},\quad (\xi^{(1)}\xi^{(2)})_{ik},\quad\ldots,\quad
        (\xi^{(1)}\xi^{(2)}\cdots\xi^{(n-1)})_{i_1i_2\cdots i_{n-1}}.
\]
Somit kommt als Spezialisierung von [5]:
\[
\begin{gathered}
        S=AB;\qquad
        BF=F\bigl(\xi_i^{(1)},(\xi^{(1)}\xi^{(2)})_{ik},\ldots\bigr),\\
        ABF=\left(\sum t_i\frac{\p}{\p\xi_i^{(1)}}\right)^{\lambda_1}
        \left(\sum t_{ik}\left(\frac{\p}{\p\xi^{(1)}}\frac{\p}{\p\xi^{(2)}}\right)_{ik}\right)^{\lambda_2}
        \cdots F\bigl(\xi_i^{(1)},(\xi^{(1)}\xi^{(2)})_{ik},\ldots\bigr),
\end{gathered}
\]
wo $\lambda_1,\lambda_2,\ldots$ die Gradzahlen in $t_i,t_{ik},\ldots$ bedeuten.

Aus invariantentheoretischen Überlegungen (es lassen sich nämlich alle Invarianten von $F$, die linear in den Koeffizienten sind, linear zusammensetzen aus $\Phi$ und aus Formen aus $M$, also insbesondere auch das nicht zu $M$ gehörende $SF$)\footnote{Dies folgt etwa aus \S~6 und \S~7, 2. meiner zitierten Arbeit ,,Zur Invariantentheorie der Formen von $n$ Variabeln``.} folgt hier die Kongruenz: $\Phi\equiv a_1SF\modu{M}$, und wegen der Normalformeigenschaft daraus anstelle von [4] die einfache Relation:
\begin{equation}
        \Phi=a_1SF=a_1S\Phi.
\tag{10}
\end{equation}
Die zweite Gleichung gilt, da die Anwendung von $S$ jede Form aus $M$ zum Verschwinden bringt. Formel [10] zeigt, daß die Normalform $\Phi$ selbst zugleich Eigenfunktion ist; und zwar gibt es nach der obigen Invariantenüberlegung keine weiteren Eigenfunktionen, die zugleich Invarianten von $F$ sind (vgl. zu [10] Deruyts, Formel (10) und (10')).

\begin{center}
\textbf{Berichtigungen.}
\end{center}

1. In der Arbeit ,,Die allgemeinsten Bereiche aus ganzen transzendenten Zahlen`` (Math. Ann. 77) ist Seite 121, Zeile 8 von oben, Z. 9 v. o. und Z. 11 v. u. der Körper $(H,Y)$ zu ersetzen durch den ,,Körper $(H,Y,Y',\ldots,Y^{(\lambda_0-1)})$, wo $Y',\ldots,Y^{(\lambda_0-1)}$ die Konjugierten von $Y$ inbezug auf $H$ bedeuten.`` Entsprechend ist S. 120, Z. 5 v. o. u. S. 121, Z. 9 von o. $(H,y)$ zu ersetzen durch $(H,y,y',\ldots,y^{(\lambda_0-1)})$.

2. In der Arbeit ,,Gleichungen mit vorgeschriebener Gruppe`` (Math. Ann. 78) muß es bei der Formulierung des Hilbertschen Irreduzibilitätssatzes -- S. 222, Z. 3 v. u. -- anstatt ,,Zahlkörper $\Omega$`` genauer heißen: ,,endlicher algebraischer Zahlkörper $\Omega$``, worauf mich Herr Ostrowski aufmerksam macht. In der Tat ist ja etwa in bezug auf den Körper aller algebraischen Zahlen die Gruppe jeder Gleichung mit algebraischen Zahlkoeffizienten die identische. -- Entsprechend ist in der Einleitung unter ,,beliebiger Rationalitätsbereich`` zu verstehen ein endlicher algebraischer Zahlkörper, dem noch endlich viele Parameter adjungiert werden können, oder einer der so entstehenden Zwischenkörper.

\begin{center}
(Angenommen September 1919.)
\end{center}

\end{document}
